problem:  
Let \(M^n\) be a compact, simply connected Riemannian manifold with positive sectional curvature, and let \(G = \mathrm{Isom}_0(M)\) denote the identity component of its isometry group.  Assume that  
\[
\dim(G) \ge c n
\]
for some fixed constant \(c > 0\).  Known positive‑curvature models fit two extremes: either \(G\) acts transitively (the rank‑one symmetric spaces, with \(\dim(G) = O(n^2)\)) or \(G\) acts with very limited symmetry (Eschenburg, Aloff–Wallach, or Bazaikin spaces, with small \(\dim(G)\)).  

**Question:**  
Does there exist a universal constant \(c_0 < 1\) such that, for all sufficiently large \(n\), the following implication holds: if \(M^n\) has positive sectional curvature and \(\dim(G) \ge c_0 n\), then \(M\) is diffeomorphic to a rank‑one symmetric space?  
Equivalently, can one prove that any positively curved manifold whose isometry group has linearly large dimension (but still far below the maximal \(O(n^2)\) scale) must already exhibit the same global topological rigidity as the rank‑one symmetric models?

Why is it a "good" problem:  
This problem probes the boundary between **quantitative symmetry** and **geometric rigidity**.  It asks whether there exists a universal proportional threshold in isometry‑group dimension that enforces rank‑one topology in positive curvature, extending traditional symmetry‑rank results into a new regime.  It calls for combining insights from **Lie group theory** (dimension and rank growth), **transformation group geometry** (orbit‑type stratification), and **global Riemannian geometry** (curvature and fixed‑point structure).  

If true, such a result would reveal a *phase transition* in the geometry of positively curved manifolds: once the continuous symmetry dimension exceeds a linear fraction of the manifold’s dimension, curvature positivity alone forces the geometry to be one of the known rank‑one models.  If false, any counterexample would uncover a new family of highly symmetric positively curved manifolds—currently unknown.  The problem’s formulation is elementary yet its resolution demands broad, deep insight across multiple mathematical domains, making it a genuinely good and far‑reaching research question.